% this file is called up by thesis.tex
% content in this file will be fed into the main document

\chapter{Implementation Details} \label{chap:Implementation Details} % top level followed by section, subsection


% ----------------------- paths to graphics ------------------------

% change according to folder and file names
\ifpdf
    \graphicspath{{4_chapter4/figures/PNG/}{4_chapter4/figures/PDF/}{4_chapter4/figures/}}
\else
    \graphicspath{{4_chapter4/figures/EPS/}{4_chapter4/figures/}}
\fi


% ----------------------- contents from here ------------------------

\section{Data Preprocessing}

The primary aim of preprocessing was to convert all features, regardless of their original datatype (numeric, categorical, text based, or mixed formats), into a standardised numerical representation suitable for machine learning. All non response and sentinel values were consistently mapped to $\texttt{-1.0}$, and all valid values were converted to $\texttt{float64}$. This ensured compatibility across the five \gls{TILDA} waves and enabled uniform downstream feature engineering and modelling.
\subsection{Data Overview}
% explain this is what we have after teh mnual cleaing described in chapter 3 we used this to visuaise in depth etc etc and tehne went on to clean...
\begin{table}[htbp]
\centering
\caption{Dimensions (Rows and Features) of Raw and Harmonised Datasets}
\label{tab:dataset_shape}
\begin{tabular}{|c|c|c|}
\hline
\textbf{Dataset} & \textbf{Row (participants)} & \textbf{Column (variables)} \\
\hline
Wave 1 & 8501 & 809 \\
Wave 2 & 7206 & 482 \\
Wave 3 & 6397 & 639 \\
Wave 4 & 5713 & 592 \\
Wave 5 & 4978 & 577 \\
Harmonised & 8501 & 453 \\
\hline
\end{tabular}
\end{table}

\subsection{Cleaning Engine}

The preprocessing strategy was built around a systematic classification of raw variable types. An exploratory audit across the five waves revealed nine distinct data scenarios (summarised in Appendix~\ref{tab:cleaning_scenario_tables}), ranging from simple numeric fields to complex multi-class string hybrids. The challenge was not merely to clean variables, but to construct a reproducible pipeline that produced \emph{numerical homogeneity} across all waves.

To accomplish this, the nine scenarios were organised into a \emph{Six-Step Cleaning Engine}. The steps reflect a principle of cumulative simplification: resolving simple and unambiguous cases first increases the proportion of clean numerical variables before addressing more complex structures.

\begin{itemize}
    \item \textbf{Steps 1–2: Missingness and sentinel standardisation.} These steps ensured consistent encoding of \texttt{-1.0} across all waves.
    \item \textbf{Step 3: Numeric recovery.} Units, symbols, and censored numeric strings (e.g., ``50+'') were stripped and converted.
    \item \textbf{Steps 4–6: Categorical complexity.} Ordinal, binary, and multi-class textual variables were processed in order of transformation confidence, with Step 6 applying the most general parsing rules.
\end{itemize}

This structured sequence ensured that each wave passed through the same deterministic cleaning logic. Figure~\ref{fig:cleaning_flow} shows the flow of the pipeline.

% --- FIGURE ---
\begin{figure}[htbp]
\centering
\begin{tikzpicture}[
    node distance=12mm,
    every node/.style={rectangle, rounded corners, draw=black, align=center, minimum width=5.2cm, minimum height=1.1cm}
]

\node (step1) {Step 1:\\Empty / Blank Handling};
\node (step2) [below=of step1] {Step 2:\\Sentinel Standardisation};
\node (step3) [below=of step2] {Step 3:\\Numeric Cleaning \& Unit Removal};
\node (step4) [below=of step3] {Step 4:\\Censored / Ordinal Mapping};
\node (step5) [below=of step4] {Step 5:\\Binary Classification};
\node (step6) [below=of step5] {Step 6:\\Multi-class / Mixed Parsing};

\draw[->, thick] (step1) -- (step2);
\draw[->, thick] (step2) -- (step3);
\draw[->, thick] (step3) -- (step4);
\draw[->, thick] (step4) -- (step5);
\draw[->, thick] (step5) -- (step6);

\end{tikzpicture}
\caption{Flowchart of the Six-Step Cleaning Pipeline applied to all \gls{TILDA} waves. See count tables in Appendix~\ref{tab:cleaning_scenario_tables} for detailed scenario progression at each step.}
\label{fig:cleaning_flow}
\end{figure}

\begin{enumerate}
    \item \textbf{Empty / Blank Values:} Converts all \texttt{NaN}, empty strings, and whitespace to \texttt{-1.0}.
    \item \textbf{Sentinel Values:} Standardises wave specific non-response codes (e.g., $-98, -99$) to \texttt{-1.0}.
    \item \textbf{Numeric Transformation:} Removes non-numeric characters and converts censored or annotated values to numerical equivalents.
    \item \textbf{Censored Categorical:} Applies known ordinal transformations (e.g., health scales, Likert ranges).
    \item \textbf{Binary Categorical:} Automatically maps two level variables (e.g., Yes/No, Male/Female) to \texttt{0.0/1.0}.
    \item \textbf{Multi-Class / Mixed:} Applies text parsing, numerical extraction, and rule based mappings to remaining multi-class variables.
\end{enumerate}

\section{Data Integration and Longitudinal Alignment}

Following cleaning, features were aligned across waves based on:

\section{Summary}

This chapter detailed the preprocessing pipeline applied to the five \gls{TILDA} waves. The approach centred on a reproducible Six-Step Cleaning Engine, combined with systematic wave-to-wave standardisation and thematic feature organisation. These steps produced a structured dataset suitable for longitudinal analysis and machine learning modelling. No cross-cohort harmonisation was performed; any references to a “Harmonised” dataset refer solely to the external RAND Harmonised TILDA.





% ---------------------------------------------------------------------------
% ----------------------- end of thesis sub-document ------------------------
% ---------------------------------------------------------------------------